\documentclass[11pt,a4paper]{article}

\usepackage[margin=1in]{geometry}
\usepackage{times}
\usepackage{booktabs}
\usepackage{longtable}
\usepackage{array}
\usepackage{xcolor}
\usepackage{enumitem}
\usepackage[hidelinks]{hyperref}
\usepackage{titlesec}

\definecolor{hdr}{RGB}{234,240,246}
\newcommand{\code}[1]{\texttt{\detokenize{#1}}}
\newcommand{\yes}{\textbf{Yes}}
\newcommand{\no}{No}
\newcommand{\partialy}{Partial}

\titleformat{\section}{\large\bfseries}{\thesection}{0.6em}{}
\titleformat{\subsection}{\normalsize\bfseries}{\thesubsection}{0.6em}{}

\setlist[itemize]{itemsep=1pt,topsep=3pt,leftmargin=1.3em}

\title{\textbf{PTV VISSIM and DhakaSim (Python): A Comparative Assessment}\\[4pt]
\large Capabilities, Distinctive Strengths, and a Path to Professional Parity}
\author{Mehzad Galib}
\date{\today}

\begin{document}
\maketitle

\begin{abstract}
\noindent
PTV VISSIM is the microscopic traffic simulator most widely used by government and
consulting organisations in Bangladesh. DhakaSim is a research simulator built
specifically for the non-lane-based, heterogeneous traffic of cities such as Dhaka,
recently reimplemented in Python and extended with real survey data. This report
compares the two: what each does well, where they genuinely differ in modelling
philosophy rather than merely in polish, and what DhakaSim would need in order to
support the professional workflow that VISSIM users expect. The comparison of
professional practice follows the workflow set out in Dr.~Md.~Hadiuzzaman's
lecture \emph{Microscopic Simulation Study} (BUET). The conclusion is that the two
tools are complementary rather than competing: VISSIM is the stronger engineering
platform, while DhakaSim models a traffic regime VISSIM does not represent. The
recommended direction is therefore not to imitate VISSIM, but to add the workflow
and evaluation machinery that would make DhakaSim's distinctive model usable and
credible in professional studies.
\end{abstract}

\section{Scope and Basis of Comparison}
This document compares VISSIM against DhakaSim \emph{as currently developed} --- the
Python reimplementation together with the extensions described in
Section~\ref{sec:dhakasim}. Statements about professional simulation practice follow
the lecture's workflow: project definition, data collection, network coding, error
checking, calibration and validation, scenario design, simulation runs, result
analysis, and documentation. No code was modified in preparing this report; it is an
assessment, not a change.

\section{The Professional Workflow VISSIM Supports}
The lecture frames a simulation study as a disciplined engineering process, and
VISSIM is built around it:

\begin{itemize}
  \item \textbf{Data preparation:} geometry from construction drawings, aerial
        photos and grades; controls from the operating agency's signal timings;
        existing demand as turning volumes, OD tables and vehicle mix; calibration
        data such as speeds, queues and travel times; and forecast future demand.
  \item \textbf{Network coding:} importing and scaling an aerial or CAD overlay,
        using coding templates, laying links and nodes over the image, then coding
        link attributes (lanes, free-flow speeds), intersection attributes (control
        type and parameters, turn-lane designations, stop bars, turn pockets),
        source/sink zones, vehicle types and OD tables.
  \item \textbf{Error checking:} verifying connectivity, geometry against the aerial
        photo, controls, prohibited turns and lane restrictions, zone volumes against
        counts, and vehicle-mix proportions --- assisted by practical strategies such
        as colour-coding links by attribute, loading 50\% demand to expose coding
        errors, following a single vehicle to spot unrealistic braking or lane
        changes, and testing simple hand-computable cases.
  \item \textbf{Calibration and validation:} treated as two distinct problems ---
        \emph{supply side} (capacity and saturation flow; global car-following,
        lane-changing and gap-acceptance parameters; then link-specific fine tuning)
        and \emph{demand side} (route choice, driver familiarity, link cost
        functions, free-flow speed) --- with explicit calibration targets, validation
        against the fundamental diagram and field observation, and search algorithms
        ranging from golden-section for a single parameter to contour plots and
        genetic algorithms for two or more.
  \item \textbf{Result analysis:} animation output (snapshots, hotspots, dynamic link
        colouring, vehicle traces) alongside numerical output (volumes, travel time,
        speed, delay, stops, density, queues), summarised through system indicators
        such as VMT and VHT, localised indicators such as queue-overflow and hotspot
        reports, and proper statistics including means, variances and hypothesis
        testing.
  \item \textbf{Bias correction:} discarding a warm-up period, keeping congestion
        inside the study area and simulation period, and counting blocked or
        unreleased vehicles.
\end{itemize}

\section{DhakaSim as Currently Developed}
\label{sec:dhakasim}
DhakaSim models each road as a set of thin lateral \emph{strips} rather than lanes; a
vehicle occupies several strips and may move to any free strip. This is the core
design decision, and it is what makes the tool able to represent traffic that has no
lane discipline. The present build provides:

\begin{itemize}
  \item a validated Python reimplementation of the original Java simulator, with a
        numeric-compatibility layer reproducing Java arithmetic exactly, pinned by a
        regression test;
  \item thirteen vehicle types spanning motorised and non-motorised traffic, and
        side-friction modelling: parked cars, rickshaws and CNGs, standing
        pedestrians, pedestrians walking along the road, and road-crossing
        pedestrians;
  \item a library of thirteen car-following models and four lane-changing models,
        selectable per run;
  \item five real intersections built from 24-hour classified turning-count surveys
        --- the Kakrail corridor (two junctions), Bijoy Sarani, Khamarbari, Banani~23
        and Banani~27 --- each with hourly demand and an hourly vehicle mix, selectable
        from the GUI along with the hour of day;
  \item an animated GUI with pan and zoom, per-type vehicle colours, named nodes, a
        collapsible legend and a live per-type vehicle counter;
  \item an automatic per-run HTML report containing headline metrics, a configuration
        table, per-type results, charts, an embedded run animation and a glossary,
        plus raw CSV outputs.
\end{itemize}

\section{Head-to-Head Comparison}

\begin{longtable}{@{}p{0.26\textwidth}p{0.34\textwidth}p{0.34\textwidth}@{}}
\toprule
\textbf{Aspect} & \textbf{PTV VISSIM} & \textbf{DhakaSim (Python)} \\
\midrule
\endhead
Traffic regime &
Lane-based; heterogeneity supported but within lane discipline &
\textbf{Non-lane-based by construction} --- strip occupancy, free lateral movement \\
\addlinespace
Vehicle population &
User-defined classes, typically motorised &
13 types including rickshaw, van/cart, CNG, bicycle; mix measured per site and hour \\
\addlinespace
Side friction &
Not represented as a first-class phenomenon &
\textbf{Explicitly modelled}: parked vehicles, standing/walking/crossing pedestrians \\
\addlinespace
Car-following &
Wiedemann 74/99 psycho-physical models &
\textbf{13 selectable models}, including a strip-based weighted model \\
\addlinespace
Lane changing &
Rule-based with gap acceptance &
4 selectable models (naive, GHR, Gipps, MOBIL) \\
\addlinespace
Network coding &
Graphical editor over scaled aerial/CAD background; templates; turn pockets, stop bars &
Hand-edited text files; no editor, no background image \\
\addlinespace
Demand input &
OD matrices, turning volumes, dynamic assignment, time-sliced demand &
OD demand per hour from survey counts; fixed shortest-path routes, constant rate within a run \\
\addlinespace
Signal control &
Fixed-time, actuated, VAP/RBC logic, priority, external controllers &
Fixed round-robin phase change only (adaptive control present but disabled) \\
\addlinespace
Calibration support &
Documented supply/demand strategy; parameter sets; multi-run seeds &
Manual; single legacy corridor calibration; seed handling prevents repeatability \\
\addlinespace
MOEs &
Volumes, speed, delay, stops, density, queues, travel time, VMT/VHT, node evaluation &
Speed, waiting time, trip time, fuel, flow, per-type and per-route CSVs; \textbf{no queue, density, delay, VMT/VHT} \\
\addlinespace
Statistical rigour &
Multi-run replication, warm-up handling, confidence intervals &
Single run; no warm-up exclusion, no replication or interval estimates \\
\addlinespace
Visual analysis &
2D/3D animation, dynamic link colouring, hotspots, vehicle traces, video export &
2D animation, per-type colours, live counter, embedded SVG run animation \\
\addlinespace
Scenario management &
Scenario manager, alternatives, comparison reports &
One network and hour per run; no scenario comparison \\
\addlinespace
Extensibility &
COM API, external driver models, scripting &
\textbf{Open Python source}, directly modifiable and scriptable \\
\addlinespace
Cost / access &
Commercial licence &
\textbf{Free, no dependencies}, standard-library Python \\
\bottomrule
\end{longtable}

\section{What Makes DhakaSim Distinctive}
Three of DhakaSim's properties are not matters of maturity, and would not be obtained
by adopting VISSIM:

\begin{enumerate}
  \item \textbf{It models a traffic regime VISSIM does not.} VISSIM assumes lane
        discipline as its organising principle. Where vehicles of very different
        sizes and speeds share an undivided carriageway and position themselves
        freely across it, the lane abstraction misrepresents the mechanism of
        congestion, not merely its magnitude. DhakaSim's strips represent that
        behaviour directly.
  \item \textbf{Side friction is a modelled phenomenon, not a fudge factor.} In
        professional practice the effects of parked vehicles and pedestrian
        interference are usually absorbed into capacity adjustment factors.
        DhakaSim generates those elements explicitly, which is what allows their
        contribution to be studied rather than assumed.
  \item \textbf{It is open, inspectable and free.} Every behavioural rule can be
        read and changed; thirteen car-following models can be swapped and compared;
        there is no licence barrier for a university or a public agency. This is
        also what makes it a viable research platform for machine-learning work,
        which a closed commercial tool is not.
\end{enumerate}

\section{Where VISSIM Is Ahead}
The gaps are concentrated in engineering workflow rather than in traffic physics:
network construction is manual and error-prone without a graphical editor or a scaled
background image; signal control is limited to fixed round-robin phasing, which rules
out the signal-design studies that motivate many professional projects; the standard
intersection performance measures --- queue length, delay, density, stops, VMT/VHT ---
are absent; and there is no support for the statistical hygiene the lecture treats as
mandatory, namely warm-up exclusion, multi-run replication and confidence intervals.
The last point is aggravated by a known defect: the random seed in the parameter file
is ignored, so headless runs are not reproducible.

\section{Improvement Roadmap}
The following recommendations are ordered by the ratio of professional value to
implementation effort. They are intentionally framed so that none of them compromises
the model's validated numerical equivalence with the original Java implementation ---
any change to vehicle movement rules should be made deliberately and documented, not
as a side effect of adding features.

\subsection*{Tier 1 --- Credibility (highest value, modest effort)}
\begin{enumerate}
  \item \textbf{Honour the random seed and add multi-run replication.} Make runs
        reproducible, then add an $n$-replication mode that reports the mean and a
        confidence interval per measure. Without this, no calibration result and no
        scenario comparison is defensible.
  \item \textbf{Add warm-up handling.} Exclude an initial loading period from the
        statistics, as the lecture requires, and report the excluded interval.
  \item \textbf{Add the missing MOEs:} queue length and queue-overflow reporting,
        control delay, stops, density, and network totals VMT and VHT. These are the
        measures agencies expect, and most are derivable from state the simulator
        already tracks.
  \item \textbf{Repair the accident and collision counters}, which currently report
        zero while events are written only to the log.
\end{enumerate}

\subsection*{Tier 2 --- Workflow (high value, moderate effort)}
\begin{enumerate}\setcounter{enumi}{4}
  \item \textbf{Signal control worthy of the name:} per-approach phases with green
        splits, cycle length and offsets, defined in an input file; then re-enable
        and validate the existing adaptive controller. This unlocks the signal
        studies that dominate professional practice, and connects directly to the
        group's NSGA-II signal-scheduling research.
  \item \textbf{Scaled background imagery and a network editor.} Even a read-only
        georeferenced background image behind the animation would let networks be
        checked against reality; a simple graphical editor for nodes and links would
        remove the main barrier to building new sites. The four synthetic star
        geometries currently in use are the clearest symptom of this gap.
  \item \textbf{Error-checking tools} following the lecture's checklist: connectivity
        and demand/route consistency checks, a reduced-demand diagnostic mode, and
        single-vehicle tracing to expose implausible braking or strip changes.
  \item \textbf{Time-varying demand within a run.} Demand is currently constant for
        the whole run, so an hour is simulated as steady conditions. Supporting
        demand profiles that change during a run would permit peak-spreading studies.
\end{enumerate}

\subsection*{Tier 3 --- Analysis and Presentation (valuable, larger effort)}
\begin{enumerate}\setcounter{enumi}{8}
  \item \textbf{Scenario management:} named scenarios with automatic side-by-side
        comparison in the HTML report --- the natural extension of the existing
        per-run reporting, and the format agencies actually consume.
  \item \textbf{Richer visual diagnostics:} dynamic link colouring by speed or
        density, hotspot highlighting, and vehicle trajectory traces, all of which
        the lecture lists as standard animation output.
  \item \textbf{Validation against the fundamental diagram:} flow--density--speed
        plots from simulated data compared with field observation, which is the
        lecture's stated validation criterion and is currently absent.
  \item \textbf{Automated calibration.} The lecture's own progression --- golden
        section, contour plots, genetic algorithms --- is the conventional route.
        Given the simulator's runtime, a surrogate-assisted approach is preferable,
        and this aligns with the planned MSc work on machine-learning-based
        calibration.
\end{enumerate}

\subsection*{Deliberately Not Recommended}
Pursuing 3D visualisation, a COM-style external API, or full dynamic traffic
assignment would consume disproportionate effort for little research return, and
would move DhakaSim toward competing with VISSIM on VISSIM's own terms. The
comparative advantage lies elsewhere.

\section{Priority Summary}

\begin{table}[h]
\centering
\small
\begin{tabular}{@{}p{0.40\textwidth}ccp{0.28\textwidth}@{}}
\toprule
\textbf{Improvement} & \textbf{Value} & \textbf{Effort} & \textbf{Unlocks} \\
\midrule
Seed fix + replication + warm-up & High & Low & Defensible results \\
Queue, delay, density, VMT/VHT   & High & Low--Med & Agency-standard reporting \\
Phase-based signal control        & High & Medium & Signal design studies \\
Background image / network editor & High & Medium--High & Realistic new sites \\
Error-checking toolkit            & Med  & Low & Trustworthy coding \\
Time-varying demand               & Med  & Medium & Peak spreading \\
Scenario comparison               & Med  & Medium & Alternatives analysis \\
Fundamental-diagram validation    & High & Medium & Model credibility \\
Automated calibration             & High & High & Thesis contribution \\
\bottomrule
\end{tabular}
\caption{Improvement priorities by value and effort.}
\end{table}

\section{Conclusion}
VISSIM is the more complete engineering platform, and for a conventional
lane-disciplined study in Bangladesh it remains the appropriate tool. DhakaSim is not
a weaker substitute for it; it is a different instrument, one that represents
non-lane-based heterogeneous traffic and its side-friction elements in a way VISSIM
structurally cannot, and it is open and free where VISSIM is closed and licensed.

The most productive direction is therefore not to replicate VISSIM's feature list but
to close the workflow gap around DhakaSim's distinctive model: make results
reproducible and statistically sound, report the performance measures practitioners
expect, give it credible signal control, and reduce the cost of building an accurate
network. Those changes would let DhakaSim's model be taken seriously in professional
studies, which is precisely where the value of a Dhaka-specific simulator lies.

\vspace{1em}
\noindent\textit{Reference:} Md.~Hadiuzzaman, \emph{Lecture 4: Microscopic Simulation
Study}, Department of Civil Engineering, BUET.

\end{document}
